% Table S22 — Polarity-inversion contrast: grounded vs learned readout
\begin{table}[t]
\centering
\setlength{\tabcolsep}{6pt}
\caption{%
Effect of prototype polarity inversion on the grounded routing readout versus the unconstrained learned classification head of the original ERP-XTTN \cite{wyman2024}, at the three-channel montage (mean AUROC, LOSO, seed-averaged across five seeds). Both readouts operate over the identical attention tensor and identical frozen prototypes; only the readout differs. Inversion is applied at inference on the frozen model. A readout that depends on prototype content should reverse under inversion; one that classifies from attention-pattern geometry should be unaffected. Noise-null values for the grounded arm are in Table~\ref{sup:tab-grounding-3-routing}. Datasets ordered by mean AUROC across all five methods, descending.%
}
\label{sup:tab-polarity-contrast}
\begin{tabular}{lcccc}
\toprule
Dataset & Grounded intact & Grounded polarity-inv. & Learned head intact & Learned head polarity-inv. \\
\midrule
HRI & 0.81 $\pm$ 0.10 & 0.19 $\pm$ 0.10 & 0.82 $\pm$ 0.11 & 0.82 $\pm$ 0.11 \\
ERN & 0.81 $\pm$ 0.11 & 0.19 $\pm$ 0.11 & 0.82 $\pm$ 0.12 & 0.82 $\pm$ 0.12 \\
P300 & 0.64 $\pm$ 0.07 & 0.36 $\pm$ 0.07 & 0.63 $\pm$ 0.08 & 0.63 $\pm$ 0.08 \\
N400 & 0.54 $\pm$ 0.07 & 0.46 $\pm$ 0.07 & 0.51 $\pm$ 0.06 & 0.51 $\pm$ 0.06 \\
\bottomrule
\end{tabular}
\end{table}
